\section{Conclusion}

We developed a complete topology-aware cooperative caching framework for IoT edge networks. The study produced a formal problem statement, research questions, dataset strategy, executable pipeline, implemented baselines, DQN refinement, adaptive policy selection, and measured results.

The results show that topology-aware and adaptive strategies outperform simple local and global caching baselines. DQN refinement improves a greedy warm start under nominal and high-churn conditions, while diversity-aware placement is strongest under moderate perturbation. Adaptive TACC combines these behaviors and achieves the best or tied-best objective across all tested perturbation rates.

The central lesson is that cooperative edge caching should not be designed as a single static placement rule. Mobility-derived topology, perturbation level, cache diversity, and relocation stability all matter. A robust system should therefore combine structural optimization, bounded learning-based refinement, and adaptive selection across operating regimes.

Future work should evaluate the same pipeline on the raw Dartmouth/IEEE DataPort archive, add temporal train-test splits, sweep content skew and cache capacity, compare against online LRU/LFU traces, and replace explicit graph features with a graph neural policy for larger deployments.

The final artifacts include both the research manuscript and the executable components needed to reproduce the reported tables. The resulting work connects research questions, implemented methods, measured results, and limitations in a single reproducible workflow.
